\chapter*{Epilogue: The Evolving Protocol}
\addcontentsline{toc}{chapter}{Epilogue: The Evolving Protocol}
\markboth{Epilogue}{Epilogue}
\label{ch:epilogue}

\epigraph{``I'll be right here.''}{E.T., \emph{E.T. the Extra-Terrestrial}
(1982)}

Every book about a moving target ends the same way, with the author waving
vaguely at the future and hoping. I want to do slightly better than that, so this
epilogue is about which parts of what you just read will survive, and why I think
so.

\section{Reading the trajectory}

Three bets are visible in \texttt{2026-07-28}, and they tell you where this is
going.

\subsection{Statelessness as an infrastructure bet}

Deleting the handshake was a decision that MCP servers should be ordinary web
services: deployable on the same infrastructure, scaled by the same tools,
operated by the same people.

Everything followed from it. No sessions, no resumability, no server-initiated requests, MRTR in place of callbacks, explicit handles in
place of implicit state.
Five deletions and one addition, all serving one idea.

That bet is unlikely to be reversed. It is far easier to add a stateful extension
later than to remove statefulness from a core, and the whole point of the
extensions framework is that the stateful things now have somewhere to live.

\subsection{Extensions as the growth mechanism}

The core got \emph{smaller} in this revision. Tasks moved out. That almost never
happens to a protocol, and it happened because there was somewhere for it to go.

Tasks and Apps are the template. Both are official, both negotiated through
capabilities, both versioned on their own schedule. Expect more: durable
subscriptions, richer authorization, capability attestation. Expect them as
extensions, and expect the core to stay boring.

Boring is the correct ambition for a core protocol. HTTP/1.1's core is boring.
TCP's core is boring. The interesting parts live above.

\subsection{The registry as distribution}

The least developed of the three and probably the most consequential.

A protocol without distribution is a specification. A protocol with a registry is
an ecosystem, and the shape of that registry decides who gets to participate and
on what terms. Chapter~\ref{ch:access-control} treated it as a governance
artefact, which is what it is: the place where ownership, data classification,
and approval records live.

Watch this one. The technical decisions are mostly made; the distribution
decisions are not.

\section{Cross-protocol convergence}

The question people ask most: how does this fit with everything else?

\begin{table}[htbp]
\centering
\small
\begin{tabularx}{\textwidth}{@{}lXX@{}}
\toprule
\thead{Protocol} & \thead{Connects} & \thead{Boundary} \\
\midrule
MCP & An agent to capabilities & Tools, resources, prompts \\
A2A & An agent to other agents & Task delegation between peers \\
Provider function calling & One model to one app's functions & In-process \\
\bottomrule
\end{tabularx}
\caption{Three different problems. The middle column is the whole distinction.}
\label{tab:convergence}
\end{table}

My read is that these settle where they are. MCP owns the agent-to-capability
boundary. A2A owns agent-to-agent. Provider-native function calling owns the
single-application case where no boundary exists and none is needed.

Chapter~\ref{ch:multi-agent} showed the composition: an agent that consumes MCP
tools on one side and exposes MCP tools on the other is an ordinary thing, and it
needed no extension. Whether the peer-to-peer case ends up as A2A or as MCP with
another extension is, honestly, a coin flip, and it will be decided by which ecosystem gets there first.

\section{Open problems}

Three things nobody has solved, in rough order of how much they will hurt.

\subsection{Prompt injection}

Chapter~\ref{ch:threats} was pessimistic on purpose. There is no known reliable
defence, and everything in that chapter was about limiting blast radius rather
than preventing the attack.

The honest position is that this is a property of the technology. Language models process instructions and data in
one stream. Until that changes, or until something at the model layer provides a
real boundary, the trifecta framing is the best tool available, and it bounds the damage
without preventing the attack.

If someone solves this properly, a good third of Chapter~\ref{ch:threats} becomes
obsolete and I will be delighted.

\subsection{Standardised evals}

Chapter~\ref{ch:measuring} argued that evals are the only way to know whether
your system works. It could not tell you which evals, because there is no shared
suite, no shared scoring, and no way to compare your tool-selection accuracy with
anybody else's.

Which means every team builds their own, and none of the results transfer. The
protocol has a conformance-test story for wire behaviour. It has nothing for
behavioural quality, and that gap is where most of the actual difficulty lives.

\subsection{Capability attestation}

Chapter~\ref{ch:architecture} was blunt: \texttt{serverInfo} is self-reported and
unverified, and you must not make security decisions with it.

That is correct and unsatisfying. In a registry-mediated ecosystem you want to
know that this server is the one that was reviewed, that its tool definitions are
the ones that were approved, and that the App template is the one that was
scanned. Chapter~\ref{ch:threats} suggested pinning hashes, which works and is
entirely local. A shared attestation mechanism would be better, and does not
exist.

\section{What will still be true}

The part I am willing to be held to.

\textbf{Round trips will still be the enemy.} It is arithmetic, not fashion.
Every extra loop iteration costs a full model turn, and that ratio holds whatever
the wire format is.

\textbf{Context will still be the scarce resource.} Windows get bigger; so do
catalogues, and so do the results people stuff into them. The discipline in
Chapter~\ref{ch:agent-loop} outlives any particular limit.

\textbf{The host will still be the trust boundary.} It follows from the architecture itself. Any design where the model is the
security boundary is broken, and will still be broken when the models are
better.

\textbf{Measurement will still beat intuition.} Chapter~\ref{ch:optimizing}
included a section on what did not work, and that section is the most durable
thing in the chapter. The specific numbers will age. The habit will not.

\textbf{Deprecated features will still be deprecated.} Twelve months minimum,
published in a registry. That promise is the reason a book can be pinned to a
revision at all.

\section{What will change}

\textbf{The version number}, obviously.

\textbf{The extensions.} Tasks and Apps are the first two. There will be more,
and some of what is an extension today will be core tomorrow.

\textbf{The SDKs.} By the time you read this they will have caught up with
\texttt{2026-07-28}, and \texttt{meridian/protocol/} will be a teaching artefact. Use the SDKs. Read the wire when something is
confusing, which will be often enough to justify having learned it.

\textbf{The dual-era bridge.} Chapter~\ref{ch:architecture} built one because
Claude Code and most 2026 clients still open with \texttt{initialize}. That is
transitional. When your clients have moved, delete
\texttt{meridian/protocol/legacy.py} and enjoy the diff. Building it as a separate wrapper was entirely so that day is a one-line
deletion.

\section{To the reader}

Here is what I would actually like you to take away.

The specific numbers in this book are from my machine on a particular Tuesday.
They will not match yours, and that is fine, because the point was never the
numbers. The point was that the numbers exist, that they came from a program you
can run, and that when somebody claims an optimisation you can ask them for
theirs.

The protocol will move. It has told you how it intends to move, which is more
than most protocols manage, and the mechanisms underneath the syntax are stable:
statelessness, capability negotiation, the trust boundary, the interception
points, the three budgets.

Build on those. Pin your version, read the changelog, keep the deprecation
registry open in a tab, and measure everything you can afford to measure.

And when your agent takes forty seconds to do something that ought to take four,
you will now know exactly where to look. That was the whole idea.

\vspace{18pt}
\noindent\emph{Krimler}\\
\emph{July 2026}
